\documentclass[aspectratio=169,10pt]{beamer}
\usetheme[numbering=fraction,progressbar=frametitle]{metropolis}
\usepackage{fontspec}
\setsansfont{Latin Modern Sans}
\usepackage{amsmath,amssymb,mathtools,graphicx,booktabs}
\usepackage{polyglossia}
\setdefaultlanguage{english}
\setbeamertemplate{navigation symbols}{}
\setbeamerfont{frametitle}{size=\large}
\setbeamerfont{normal text}{size=\normalsize}
\setlength{\parskip}{.45em}
\hypersetup{pdfauthor={},pdfsubject={Lecture slides},pdfkeywords={},pdftitle={Complex numbers and trigonometry}}
\providecommand{\N}{\mathbb{N}}
\providecommand{\Z}{\mathbb{Z}}
\providecommand{\D}{\mathbb{D}}
\providecommand{\Q}{\mathbb{Q}}
\providecommand{\R}{\mathbb{R}}
\providecommand{\C}{\mathbb{C}}
\providecommand{\vect}[1]{\overrightarrow{#1}}
\providecommand{\vi}{\vec{\imath}}
\providecommand{\vj}{\vec{\jmath}}
\providecommand{\vk}{\vec{k}}
\providecommand{\norm}[1]{\left\lVert #1\right\rVert}
\providecommand{\vc}[2]{\begin{pmatrix}#1\\#2\end{pmatrix}}
\providecommand{\deter}[2]{\det\left(#1,#2\right)}
\providecommand{\abs}[1]{\left\lvert #1\right\rvert}
\providecommand{\dd}{\,\mathrm{d}}
\providecommand{\e}{\mathrm{e}}
\providecommand{\Card}{\operatorname{Card}}
\providecommand{\E}{\operatorname{E}}
\providecommand{\V}{\operatorname{V}}
\providecommand{\degre}{^{\circ}}
\providecommand{\un}[1]{\,\mathrm{#1}}
\providecommand{\sci}[2]{#1\times10^{#2}}
\providecommand{\pgcd}{\operatorname{PGCD}}
\providecommand{\ppcm}{\operatorname{PPCM}}
\title{Complex numbers and trigonometry}
\subtitle{Lecture and worked examples}
\date{}
\author{}
\begin{document}
\begin{frame}[plain]
\vfill
{\usebeamerfont{title}\inserttitle\par}
\medskip
{\usebeamerfont{subtitle}\insertsubtitle\par}
\medskip
{\usebeamercolor[fg]{progress bar}\rule{\textwidth}{.5pt}}
\vfill
\end{frame}
\begin{frame}{Objectives}
Establish the addition and double-angle formulas; write a complex number in exponential form and choose the form suited to a computation; compute powers; use de Moivre's and Euler's formulas to transform trigonometric expressions.
\end{frame}
\begin{frame}{Introductory activity}
For every real number \(\theta\), let \(f(\theta)=\cos\theta+i\sin\theta\).
\end{frame}
\begin{frame}{Introductory activity (2)}
\begin{enumerate}
\setcounter{enumi}{0}\item Compute \(f(0)\), \(f\left(\dfrac\pi2\right)\), \(f(\pi)\). What is the modulus of \(f(\theta)\)? Where does the image of \(f(\theta)\) lie?

\item Write \(f\left(\dfrac\pi3\right)\) and \(f\left(\dfrac\pi6\right)\) in algebraic form, then compute the product \(f\left(\dfrac\pi3\right)f\left(\dfrac\pi6\right)\). Compare it with \(f\left(\dfrac\pi2\right)\).

\item Compute likewise \(f\left(\dfrac\pi4\right)^2\) and compare it with \(f\left(\dfrac\pi2\right)\).

\end{enumerate}
\end{frame}
\begin{frame}{Introductory activity (3)}
\begin{enumerate}
\setcounter{enumi}{3}\item With a calculator or a spreadsheet, compare \(f(a)f(b)\) and \(f(a+b)\) for three pairs \((a,b)\) chosen at random.

\item Which known function satisfies a relation of the same kind, \(g(a)g(b)=g(a+b)\) and \(g(0)=1\)?

\item Expand \(f(a)f(b)\). Which trigonometric formulas would one need to know to prove the conjecture?

\end{enumerate}
\end{frame}
\begin{frame}{Addition and double-angle formulas}
For all real numbers \(a\) and \(b\):
\[\cos(a-b)=\cos a\cos b+\sin a\sin b,\qquad\cos(a+b)=\cos a\cos b-\sin a\sin b,\]
\[\sin(a-b)=\sin a\cos b-\cos a\sin b,\qquad\sin(a+b)=\sin a\cos b+\cos a\sin b.\]
\end{frame}
\begin{frame}{Addition and double-angle formulas}
In a direct orthonormal frame, let \(\vect p\vc{\cos a}{\sin a}\) and \(\vect q\vc{\cos b}{\sin b}\): these vectors are unit vectors and \(\left(\vect u,\vect p\right)\equiv a\), \(\left(\vect u,\vect q\right)\equiv b\pmod{2\pi}\). By Chasles' relation, \(\left(\vect q,\vect p\right)\equiv a-b\pmod{2\pi}\). We compute \(\vect q\cdot\vect p\) in two ways:
\[\vect q\cdot\vect p=\norm{\vect q}\norm{\vect p}\cos(a-b)=\cos(a-b)\quad\text{and}\quad\vect q\cdot\vect p=\cos b\cos a+\sin b\sin a.\]
This proves the first formula. Replacing \(b\) by \(-b\) and using \(\cos(-b)=\cos b\), \(\sin(-b)=-\sin b\), we obtain \(\cos(a+b)\). Finally \(\sin(a+b)=\cos\left(\dfrac\pi2-a-b\right)=\cos\left(\left(\dfrac\pi2-a\right)-b\right)\), which the first formula expands into \(\sin a\cos b+\cos a\sin b\); we deduce \(\sin(a-b)\) by replacing \(b\) by \(-b\).
\end{frame}
\begin{frame}{Addition and double-angle formulas}
For every real number \(a\): \(\sin2a=2\sin a\cos a\) and \(\cos2a=\cos^2a-\sin^2a=2\cos^2a-1=1-2\sin^2a\).
\end{frame}
\begin{frame}{Addition and double-angle formulas}
\(\cos\dfrac{\pi}{12}=\cos\left(\dfrac\pi3-\dfrac\pi4\right)=\dfrac12\times\dfrac{\sqrt2}2+\dfrac{\sqrt3}2\times\dfrac{\sqrt2}2=\dfrac{\sqrt6+\sqrt2}4\).
\end{frame}
\begin{frame}{Product and arguments}
For all real numbers \(\theta\) and \(\theta'\): \((\cos\theta+i\sin\theta)(\cos\theta'+i\sin\theta')=\cos(\theta+\theta')+i\sin(\theta+\theta')\).

Consequently, for all nonzero complex numbers \(z\), \(z'\) and every integer \(n\):
\[\arg(zz')\equiv\arg z+\arg z',\qquad\arg\dfrac1z\equiv-\arg z,\]
\[\arg\dfrac z{z'}\equiv\arg z-\arg z',\qquad\arg\left(z^n\right)\equiv n\arg z\pmod{2\pi}.\]
\end{frame}
\begin{frame}{Product and arguments}
Expanding gives \((\cos\theta\cos\theta'-\sin\theta\sin\theta')+i(\sin\theta\cos\theta'+\cos\theta\sin\theta')\), which the addition formulas simplify. If \(z=r(\cos\theta+i\sin\theta)\) and \(z'=r'(\cos\theta'+i\sin\theta')\), then \(zz'=rr'(\cos(\theta+\theta')+i\sin(\theta+\theta'))\) with \(rr'>0\): this is a trigonometric form of \(zz'\). The other relations follow: \(z\times\dfrac1z=1\) has argument \(0\), then \(\dfrac z{z'}=z\times\dfrac1{z'}\), and an induction for \(n\ge0\), extended to \(n<0\) through the inverse.
\end{frame}
\begin{frame}{Product and arguments}
Multiplying by \(\cos\theta+i\sin\theta\) turns through an angle \(\theta\) about \(O\). The conjecture of the activity of chapter 2 is proved: multiplying by \(i\) is a direct quarter turn, multiplying by \(1+i=\sqrt2\left(\cos\dfrac\pi4+i\sin\dfrac\pi4\right)\) is an eighth of a turn followed by an enlargement with ratio \(\sqrt2\).
\end{frame}
\begin{frame}{Exponential form}
For every real number \(\theta\), we write \(\e^{i\theta}=\cos\theta+i\sin\theta\). Every nonzero complex number can be written \(z=r\e^{i\theta}\) with \(r=\abs z\) and \(\theta\equiv\arg z\pmod{2\pi}\): this is its \emph{exponential form}.
\end{frame}
\begin{frame}{Exponential form}
For all real numbers \(\theta\), \(\theta'\) and every integer \(n\):
\[\e^{i\theta}\e^{i\theta'}=\e^{i(\theta+\theta')},\qquad\dfrac1{\e^{i\theta}}=\e^{-i\theta}=\overline{\e^{i\theta}},\qquad\dfrac{\e^{i\theta}}{\e^{i\theta'}}=\e^{i(\theta-\theta')},\qquad\left(\e^{i\theta}\right)^n=\e^{in\theta}.\]
Values to know: \(\e^{i0}=1\), \(\e^{i\pi/2}=i\), \(\e^{i\pi}=-1\), \(\e^{-i\pi/2}=-i\), \(\abs{\e^{i\theta}}=1\).
\end{frame}
\begin{frame}{Exponential form}
For every real number \(\theta\) and every integer \(n\): \((\cos\theta+i\sin\theta)^n=\cos(n\theta)+i\sin(n\theta)\).
\end{frame}
\begin{frame}{Exponential form}
For every real number \(\theta\): \(\cos\theta=\dfrac{\e^{i\theta}+\e^{-i\theta}}2\) and \(\sin\theta=\dfrac{\e^{i\theta}-\e^{-i\theta}}{2i}\).
\end{frame}
\begin{frame}{Exponential form}
\(\e^{i\theta}+\e^{-i\theta}=\e^{i\theta}+\overline{\e^{i\theta}}=2\operatorname{Re}\left(\e^{i\theta}\right)=2\cos\theta\), and \(\e^{i\theta}-\e^{-i\theta}=2i\operatorname{Im}\left(\e^{i\theta}\right)=2i\sin\theta\).
\end{frame}
\begin{frame}{Exponential form}
For a sum, the algebraic form; for a product, a quotient or a power, the exponential form. Example: \(1+i\sqrt3=2\e^{i\pi/3}\), so
\[\left(1+i\sqrt3\right)^{10}=2^{10}\e^{10i\pi/3}=1{,}024\,\e^{-2i\pi/3}=-512-512\sqrt3\,i,\]
since \(\dfrac{10\pi}3-4\pi=-\dfrac{2\pi}3\).
\end{frame}
\begin{frame}{Exponential form}
Expand \((\cos\theta+i\sin\theta)^3\) with the binomial formula and identify the real part:
\(\cos3\theta=\cos^3\theta-3\cos\theta\sin^2\theta=4\cos^3\theta-3\cos\theta\).
\end{frame}
\begin{frame}{Exponential form}
\(\cos^3\theta=\dfrac{\left(\e^{i\theta}+\e^{-i\theta}\right)^3}{8}=\dfrac{\e^{3i\theta}+3\e^{i\theta}+3\e^{-i\theta}+\e^{-3i\theta}}8=\dfrac{\cos3\theta+3\cos\theta}4\).
\end{frame}
\begin{frame}{Exponential form}
To factorise a sum of exponentials, take out the exponential of the mean angle: \(\e^{ia}+\e^{ib}=\e^{i\dfrac{a+b}2}\left(\e^{i\dfrac{a-b}2}+\e^{-i\dfrac{a-b}2}\right)=2\cos\dfrac{a-b}2\,\e^{i\dfrac{a+b}2}\).
\end{frame}
\begin{frame}{Exponential form}
\(\e^{i\theta}\) is not a positive real number: its modulus is \(1\), but it can be negative (\(\e^{i\pi}=-1\)) or non-real. An expression \(r\e^{i\theta}\) is an exponential form only if \(r>0\).
\end{frame}
\end{document}
